\begin{tikzpicture}[
  x=1mm,y=1mm,
  every node/.style={font=\figurefont\footnotesize,text=BookInk},
  stage/.style={rectangle,draw=BookRule,line width=0.8pt,fill=BookPaper,text width=24mm,minimum height=13mm,align=center,inner sep=2mm},
  source/.style={rectangle,rounded corners=1.2mm,draw=FigureInput,line width=0.8pt,fill=FigureInput!8,text width=14mm,minimum height=10mm,align=center,inner sep=1.5mm},
  fusion/.style={stage,draw=FigureModel,fill=FigureModel!8,text width=20mm},
  check/.style={stage,draw=BookAmber,fill=BookAmber!8,text width=21mm},
  result/.style={stage,draw=FigureOutput,fill=FigureOutput!6,text width=15mm},
  flow/.style={knowledge arrow,draw=BookTeal,line width=0.9pt,rounded corners=1.2mm},
  control/.style={knowledge arrow,draw=BookMuted,line width=0.8pt,rounded corners=1.2mm},
  bus/.style={draw=BookMuted,line width=0.8pt},
  knowledge bus/.style={draw=BookTeal,line width=0.9pt},
  feedback bus/.style={draw=BookAmber,dashed,line width=0.8pt,rounded corners=1.2mm},
  adjust/.style={knowledge arrow,draw=BookAmber,dashed,line width=0.8pt,rounded corners=1.2mm},
  edge label/.style={font=\figurefont\scriptsize,inner sep=1pt,text=BookMuted}
]
  \node[stage] (need) at (15,12) {需求与资源\\范围、效果、预算};
  \node[stage] (select) at (47,12) {来源选择\\对象、机制、投入};

  \node[source] (human) at (83,36) {人工};
  \node[source] (rule) at (83,20) {规则};
  \node[source] (model) at (83,4) {模型};
  \node[source] (data) at (83,-12) {数据构造};
  \node[font=\figurefont\scriptsize\bfseries,text=BookTeal] at (83,46) {执行选定的获取机制};

  \node[stage,draw=FigureModel,fill=FigureModel!8] (candidate) at (126,12) {候选知识\\含来源与状态};
  \node[fusion,text width=24mm] (align) at (83,-46) {对象与语义对齐\\需要时再融合};
  \node[check] (quality) at (126,-46) {质量判断\\覆盖与可靠性\\下游效用};
  \node[result] (usable) at (159,-46) {可用结果\\或停止};

  \draw[control] (need.east) -- (select.west);
  \draw[control] (select.east) -- (67,12);
  \draw[bus] (67,-12) -- (67,36);
  \draw[control] (67,36) -- (human.west);
  \draw[control] (67,20) -- (rule.west);
  \draw[control] (67,4) -- (model.west);
  \draw[control] (67,-12) -- (data.west);

  \draw[knowledge bus] (100,-12) -- (100,36);
  \draw[flow] (human.east) -- (100,36);
  \draw[flow] (rule.east) -- (100,20);
  \draw[flow] (model.east) -- (100,4);
  \draw[flow] (data.east) -- (100,-12);
  \draw[flow] (100,12) -- (candidate.west);

  \draw[flow] ([xshift=-8mm]candidate.south) -- (118,-25) -| (align.north);
  \node[edge label,anchor=south] at (99,-24) {多份可比结果};
  \draw[flow] (align.east) -- (quality.west);
  \draw[flow] (candidate.south) -- (quality.north);
  \node[edge label,anchor=west,align=left] at (129,-16) {单一来源\\或独立目标};
  \draw[flow] (quality.east) -- (usable.west);

  \coordinate (feedback split) at (100,-64);
  \draw[feedback bus] (quality.south) |- (feedback split);
  \draw[adjust] (feedback split) -- (100,-72) -| (select.south);
  \node[edge label,text=BookAmber,anchor=south] at (70,-71) {调整来源选择};
  \draw[adjust] (feedback split) -- (64,-64) -- (64,-23) -- (67,-23) -- (67,-12);
  \node[edge label,text=BookAmber,anchor=south] at (82,-63) {调整获取机制};

  \draw[flow] (4,-82) -- (16,-82);
  \node[anchor=west,font=\figurefont\scriptsize] at (18,-82) {知识流动};
  \draw[adjust] (57,-82) -- (69,-82);
  \node[anchor=west,font=\figurefont\scriptsize] at (71,-82) {质量反馈};
\end{tikzpicture}
